\documentclass[11pt,a4paper]{article}
\usepackage[T1]{fontenc}
\usepackage[utf8]{inputenc}
\usepackage[slovene]{babel}
\usepackage{amsmath,amssymb}
\usepackage{graphicx}
\usepackage{booktabs}
\usepackage{xcolor}
\usepackage{listings}
\usepackage[margin=2.6cm]{geometry}
\usepackage[hidelinks]{hyperref}
\usepackage{placeins}

\renewcommand{\topfraction}{0.9}
\renewcommand{\bottomfraction}{0.8}
\renewcommand{\textfraction}{0.07}
\renewcommand{\floatpagefraction}{0.75}

\newcommand{\slikaali}[2]{\IfFileExists{#1}{\includegraphics[width=#2\linewidth]{#1}}%
  {\dopolni{sliko \texttt{\detokenize{#1}} ustvari \texttt{vizualizacije.py}}}}
\lstset{language=Python,basicstyle=\ttfamily\small,frame=single,breaklines=true,
        showstringspaces=false}

\title{Domača naloga 3:\\ Ničle Airyjeve funkcije z Magnusovo metodo reda 4}
\author{Urban Vesel}
\date{1. september 2026}

\begin{document}
\maketitle

\section{Opis problema}

Airyjeva funkcija $\mathrm{Ai}$ je rešitev začetnega problema
\begin{equation}
  y'' = x\,y, \qquad
  \mathrm{Ai}(0) = \frac{1}{3^{2/3}\,\Gamma(2/3)}, \qquad
  \mathrm{Ai}'(0) = -\frac{1}{3^{1/3}\,\Gamma(1/3)} .
\end{equation} 
Vrednosti $\mathrm{Ai}(0)$ in $\mathrm{Ai}'(0)$ sta zapisani v zaprti obliki
v \cite[razdelka 9.2.3 in 9.2.4]{dlmf}. Ker funkcija $\Gamma$ ni elementarna,
ju v kodi hranimo kot konstanti, zaokroženi na dvajset mest.
Naloga~\cite{skripta} zahteva, da poiščemo čim več njenih ničel (vse ležijo
na negativni polosi) na deset decimalk natančno; vgrajeni reševalci diferencialnih enačb niso
dovoljeni, \texttt{scipy.special.airy} pa služi izključno preverjanju.
Enačbo zapišemo kot sistem prvega reda $y' = A(x)\,y$ za
$y = (\mathrm{Ai}, \mathrm{Ai}')$ z
$A(x) = \bigl[\begin{smallmatrix}0&1\\ x&0\end{smallmatrix}\bigr]$
in integriramo od $0$ v levo.

\section{Metoda}

\subsection{Magnusova metoda reda 4}

Rešitev linearnega sistema na koraku je $y(x+h) = e^{\Omega(h)} y(x)$, kjer je
$\Omega$ Magnusova vrsta \cite{blanes}. Ta izraža logaritem propagatorja kot
vsoto členov, pri katerih $k$-ti prispeva $O(h^k)$; prvi je integral matrike,
drugi pa popravek, ki meri, koliko $A$ v različnih točkah komutira med seboj:
\begin{equation}
  \Omega(h) = \int_0^h A\,\mathrm d\tau
  + \frac12 \int_0^h\!\!\int_0^{\tau_1} [A(\tau_1), A(\tau_2)]\,
    \mathrm d\tau_2\,\mathrm d\tau_1 + \cdots
\end{equation}
Metoda reda~4 iz \cite{skripta} obdrži prva dva člena in ju izračuna z
dvotočkovno Gauss--Legendrovo kvadraturo, ki je točna za polinome stopnje
$\le 3$. Skrajšanje vrste bi samo po sebi dalo red~2, vendar je nastali
postopek časovno simetričen: korak $h$ in za njim korak $-h$ vrneta izhodiščno
stanje. Simetrični postopki imajo sod red, zato se red dvigne s~3 na~4.
Z $A_{1,2} = A\bigl(x + (\tfrac12 \mp \tfrac{\sqrt3}{6})h\bigr)$ je
\begin{equation}
  \sigma = \frac h2 (A_1 + A_2) - \frac{\sqrt3}{12} h^2\, [A_1, A_2],
  \qquad y^+ = e^{\sigma} y,
  \label{eq:magnus}
\end{equation}
z lokalno napako $O(h^5)$. 
Izmerjeni red je $4{,}00$ (levi graf na sliki~\ref{fig:wronskian}), kar potrjuje
zgornjo oceno.
Za naš $A$ je $[A_1, A_2] = (x_2 - x_1)
\operatorname{diag}(1, -1)$ in $x_2 - x_1 = \tfrac{\sqrt3}{3}h$, zato se
\eqref{eq:magnus} sklene v
\begin{equation}
  \sigma = \begin{bmatrix} -h^3/12 & h \\ h\,(x + h/2) & h^3/12 \end{bmatrix},
\end{equation}
komutatorja v kodi torej ni treba računati (pravilnost potrjuje test proti
splošni obliki).

\subsection{Eksponentna funkcija brezsledne $2\times 2$ matrike}

Ker je $\operatorname{tr}\sigma = 0$, da Cayley--Hamiltonov izrek
$\sigma^2 = zI$ z $z = \sigma_{11}^2 + \sigma_{12}\sigma_{21}$, in
eksponentna vrsta razpade na sodi in lihi del. Iz $\sigma^2 = zI$ namreč sledi
$\sigma^{2j} = z^j I$ in $\sigma^{2j+1} = z^j \sigma$, zato se sode potence
seštejejo v večkratnik identitete, lihe pa v večkratnik $\sigma$:
\begin{equation}
  e^{\sigma} = C(z)\,I + S(z)\,\sigma, \qquad
  C(z) = \sum_{j\ge 0} \frac{z^j}{(2j)!} = \cosh\sqrt z, \quad
  S(z) = \sum_{j\ge 0} \frac{z^j}{(2j+1)!} = \frac{\sinh\sqrt z}{\sqrt z}.
\end{equation}
$C$ in $S$ sta \emph{celi} funkciji $z$, zato formula velja za oba predznaka;
numerično uporabimo $\cosh/\sinh$ za $z > \delta$, $\cos/\sin$ za
$z < -\delta$ in nekaj členov vrste za $|z| \le \delta$. Ker je
$\det e^{\sigma} = e^{\operatorname{tr}\sigma} = 1$ \emph{eksaktno}, metoda
strukturno ohranja Wronskijan dveh rešitev, kar je natanko lastnost točnega
toka (Abel--Liouville: $W' = \operatorname{tr}A \cdot W = 0$); klasični RK4 te
lastnosti nima (slika~\ref{fig:wronskian}).

\subsection{Izbira koraka in asimptotika (WKB)}

Za $x = -\xi < 0$ enačba preide v $y'' = -\xi y$, torej v nihanje s počasi
spreminjajočo se frekvenco $\sqrt\xi$. Nastavek Liouville--Green
$y = \xi^{-1/4} e^{\pm i S(\xi)}$ s fazo $S' = \sqrt\xi$ enačbi zadosti do
relativnega $O(\xi^{-3/2})$ in da \cite[razdelek 9.7.9]{dlmf}
\begin{equation}
  \mathrm{Ai}(-\xi) \approx \pi^{-1/2}\,\xi^{-1/4}
  \cos\!\bigl(\tfrac23 \xi^{3/2} - \tfrac{\pi}{4}\bigr),
  \label{eq:wkb}
\end{equation} 
torej amplitudo $\sim |x|^{-1/4}$ in lokalno valovno dolžino $2\pi/\sqrt{|x|}$.
Korak zato izberemo kot $h = \min\!\bigl(h_0,\; 2\pi/(m\sqrt{1+|x|})\bigr)$,
kar da približno $m$ korakov na nihaj. Iz \eqref{eq:wkb} sledijo še tri
uporabne posledice. Ničle nastopijo tam, kjer je kosinus enak nič, torej pri
$\tfrac23|a_k|^{3/2} = \tfrac{\pi(4k-1)}{4}$; od tod vodilni McMahonov člen
$|a_k| = \bigl(\tfrac{3\pi(4k-1)}{8}\bigr)^{2/3}$ \cite[razdelek 9.9.6]{dlmf},
s katerim ocenimo, kako globoko moramo integrirati. Ker faza narašča z
$\tfrac23|x|^{3/2}$, je ničel do globine $x$ približno
$N(x) \approx \tfrac{2}{3\pi}|x|^{3/2}$. Odvod v ničli je
$|\mathrm{Ai}'(a_k)| \approx \pi^{-1/2}|a_k|^{1/4}$ in z globino narašča, zato
so ničle dobro pogojene in Newtonov postopek hitro konvergira.

\subsection{Iskanje in prečiščenje ničel}

Med pohodom zaznamo menjave predznaka $\mathrm{Ai}$ med sosednjima vozloma in
vsak oklep prečistimo z varovanim Newtonovim postopkom
$x \leftarrow x - \mathrm{Ai}(x)/\mathrm{Ai}'(x)$ (ob skoku iz oklepa
bisekcija). Vrednosti $(\mathrm{Ai}, \mathrm{Ai}')$ v vmesnih točkah dobimo s
kratko propagacijo od predpomnjene mreže. Zaradi dobre pogojenosti zadošča
nekaj iteracij, kar se v meritvah tudi pokaže: konvergenca je kvadratična in
oklep se v praksi ne zapusti.

\subsection{Kontrola natančnosti}

Fazna napaka metode reda 4 na nihaj je $\sim C' m^{-4}$; do globine $|a_k|$
je nihajev $\Phi/2\pi$ s $\Phi = \tfrac23 |a_k|^{3/2}$, premik ničle pa
$\delta a \approx \delta\varphi / \sqrt{|a_k|}$. Skupaj
\begin{equation}
  |\delta a_k| \approx C''\, |a_k|\, m^{-4}.
  \label{eq:zakon}
\end{equation}
Napako umerimo proti oraklju \texttt{scipy}: izmerjena konstanta $C''$ (glej
sliko~\ref{fig:napaka}) da napoved globine, do katere je pri izbranem $m$
dosežena zahtevana natančnost $10^{-10}$, kar je kvantitativen odgovor na
zahtevo po čim več ničlah.

\section{Rezultati}

Slika~\ref{fig:ai} prikazuje $\mathrm{Ai}$ na $[-12, 0]$ z izračunanimi ničlami.

\begin{figure}[htb]\centering
  \slikaali{slike/ai_nicle.png}{0.6}
  \caption{$\mathrm{Ai}$ na $[-12, 0]$ z označenimi ničlami: devet menjav
  predznaka, devet ničel, brez izpuščenih ali podvojenih.}
  \label{fig:ai}
\end{figure}

\begin{table}[htb]\centering
\caption{Prvih pet ničel v primerjavi z DLMF. Razlika je povsod manjša od
$10^{-11}$, kar je z rezervo pod zahtevanimi desetimi decimalkami.}
\begin{tabular}{rll}
\toprule
$k$ & $a_k$ (naša) & $a_k$ (DLMF) \\ \midrule
1 & $-2{,}338107410467$ & $-2{,}338107410459$ \\
2 & $-4{,}087949444137$ & $-4{,}087949444130$ \\
3 & $-5{,}520559828101$ & $-5{,}520559828096$ \\
4 & $-6{,}786708090077$ & $-6{,}786708090072$ \\
5 & $-7{,}944133587126$ & $-7{,}944133587120$ \\ \bottomrule
\end{tabular}
\end{table}

\begin{figure}[htb]\centering
  \slikaali{slike/red4.png}{0.48}\hfill
  \slikaali{slike/wronskian.png}{0.48}
  \caption{Levo: empirični red metode na $x_0=0 \to x=-1$; izmerjeni naklon
  na log-log skali je $4{,}00$, kar se natančno ujema z zahtevanim redom.
  Desno: drift Wronskijana pri koraku $h=0{,}05$ na $[-60,0]$. Magnus ostane
  ves čas na ravni strojne natančnosti (na koncu $7{,}2\cdot10^{-16}$), RK4
  pa nabere drift $3{,}6\cdot10^{-3}$, torej približno $5\cdot10^{12}$-krat
  več: to je neposreden odraz eksaktne ohranitve $\det e^{\sigma}=1$.
  Navpični skoki modre krivulje so točke kjer je drift nič, narisane so na spodnji rob $10^{-18}$.}
  \label{fig:wronskian}
\end{figure}

\begin{figure}[htb]\centering
  \slikaali{slike/napaka_nicel.png}{0.55}
  \caption{$|\Delta a_k|$ prvih 60 ničel proti \texttt{scipy.special.ai\_zeros}
  pri $m=800$, skupaj z zakonom \eqref{eq:zakon} za umerjeno konstanto
  $C'' \approx 4{,}05\cdot10^{-2}$. Vseh 60 ničel ima napako pod $10^{-10}$,
  največja izmerjena je $1{,}27\cdot10^{-11}$ pri $k=5$. 
  Izmerjena napaka z naraščajočim $k$ celo rahlo pada (od $7{,}2\cdot10^{-12}$
  pri $k=1$ na $2{,}1\cdot10^{-12}$ pri $k=60$), medtem ko zakon
  \eqref{eq:zakon} predvideva rast sorazmerno z $|a_k|$. Zakon je torej
  konzervativna zgornja meja.}
  \label{fig:napaka}
\end{figure}

V demonstraciji smo prvih 60 ničel preverili proti oraklju (slika \ref{fig:napaka}). 
Število korakov je sorazmerno s skupno fazo
$\Phi = \tfrac23 |x|^{3/2}$, prav tako pa je s $\Phi$ sorazmerno tudi število
ničel. Cena zato raste kot $|x|^{3/2}$ v globini, v številu ničel pa je
linearna. To potrjujejo meritve: 30 ničel dobimo v $159$ ms, 100 v $852$ ms in
300 v $2{,}5$ s, kar je približno enak čas na ničlo. 
Napovedanih $6800$ zanesljivih ničel v demu ne računamo do konca, ker bi to trajalo občutno dlje,
je pa napoved iz \eqref{eq:zakon} kvantitativen, preverljiv odgovor na to,
koliko ničel je pri danem $m$ še vredno računati.

\section{Primer uporabe}

\begin{minipage}{\linewidth}
\begin{lstlisting}
from dn3_airy_nicle import nicle, ai
z = nicle(30)          # prvih 30 nicel
z10 = nicle(x_min=-10) # vse nicle na [-10, 0]
Ai, dAi = ai(-3.0)      # (-0.378814, 0.314584)
\end{lstlisting}
\end{minipage}

\section{Zaključek}

Ničle Airyjeve funkcije smo poiskali z Magnusovo metodo reda 4 v zaprti
obliki, kar je pri koraku $h=0{,}05$ na $[-60,0]$ dalo Wronskijan, konstanten
do $7{,}2\cdot10^{-16}$, medtem ko je enak korak z RK4 nabral drift
$3{,}6\cdot10^{-3}$. Varovan Newton na dobro pogojenih ničlah konvergira v
nekaj korakih; prvih 60 ničel se z oraklju \texttt{scipy} ujema na manj kot
$1{,}3\cdot10^{-11}$, torej precej pod zahtevanih deset decimalk. Umerjeni
zakon $|\delta a_k| \approx C''|a_k|m^{-4}$ napove, da isti $m=800$ zanesljivo
nese do približno $6800$ ničel. Naravna izboljšava bi bila vektorizacija
pohoda (trenutno čista Python zanka) ali višji red Magnusove vrste za še
redkejše korake pri enaki natančnosti.

\begin{thebibliography}{9}
\bibitem{dlmf} \emph{NIST Digital Library of Mathematical Functions},
  \S 9.2, 9.7, 9.9, \url{https://dlmf.nist.gov/9}.
\bibitem{blanes} S.~Blanes, F.~Casas, J.~A. Oteo, J.~Ros, The Magnus
  expansion and some of its applications, \emph{Phys.\ Rep.}\ 470 (2009),
  \url{https://arxiv.org/abs/0810.5488}.
\bibitem{skripta} M.~Vuk, \emph{Numerična matematika v programskem jeziku
  Julia}, 2025, naloga~18.3.1.
\end{thebibliography}

\end{document}
